\documentclass[10pt,twocolumn]{article}

% ── Page geometry: 2-column, tight margins ──
\usepackage[a4paper,margin=1.8cm,top=1.5cm,bottom=1.5cm]{geometry}
\setlength{\columnsep}{0.6cm}

% ── Packages ──
\usepackage[utf8]{inputenc}
\usepackage{amsmath,amssymb}
\usepackage{graphicx}
\usepackage{booktabs}
\usepackage{enumitem}
\usepackage{caption}
\usepackage{xcolor}
\usepackage{tikz}
\usetikzlibrary{arrows.meta,positioning,fit,calc}
\usepackage[hidelinks]{hyperref}
\usepackage{microtype}

% ── Compact spacing ──
\usepackage{titlesec}
\titlespacing*{\section}{0pt}{8pt}{4pt}
\titlespacing*{\subsection}{0pt}{6pt}{2pt}
\titleformat{\section}{\normalsize\bfseries}{\thesection.}{0.4em}{}
\titleformat{\subsection}{\small\bfseries}{\thesubsection}{0.4em}{}
\setlength{\parskip}{2pt}
\setlength{\parindent}{0pt}
\setlist{nosep,leftmargin=1.2em}

% ── Caption style ──
\captionsetup{font=small,labelfont=bf,skip=4pt}

% ── Colors ──
\definecolor{boxgray}{HTML}{f5f4ef}

\begin{document}

% ══════════════════════════════════════════
% TITLE
% ══════════════════════════════════════════
\twocolumn[{%
\centering
{\Large\bfseries SparseTap: Sparse XOR Dependency Discovery\\via Statistical Voting and Gradient Descent\par}
\vspace{4pt}
{\normalsize KwanHyeon Park \quad$\cdot$\quad KRAFTON AI R\&D Hackathon, Round 1, Day 2\par}
\vspace{8pt}
\rule{\textwidth}{0.4pt}
\vspace{6pt}
}]

% ══════════════════════════════════════════
% 1. PROBLEM SETUP
% ══════════════════════════════════════════
\section{Problem Setup}
\label{sec:setup}

A binary sequence is generated by a hidden sparse XOR rule with additive noise:
%
\begin{equation}
s[n] = \bigoplus_{k=1}^{S} s[n - d_k] \oplus e[n], \quad e[n] \sim \mathrm{Bern}(0.2)
\label{eq:rule}
\end{equation}
%
where $d_1 < d_2 < \cdots < d_S$ are unknown offsets, $S \le 16$, and
the seed length $W = d_S \le 64$.
We are given $N{=}2000$ noisy sequences of length 256,
all sharing the same hidden offsets but with independent random seeds.
The goal is to recover the offsets and predict the next 192 bits of a
noise-free test sequence whose first 64 bits are observed.

% ══════════════════════════════════════════
% 2. APPROACH 1: STATISTICAL VOTING
% ══════════════════════════════════════════
\section{Approach 1: Majority-Vote Correlation}
\label{sec:stat}

\subsection{Key Observation}

For a candidate offset $d$, define the \emph{residual} of sequence $i$ at position $n$ as:
%
\begin{equation}
r^{(i)}_d[n] = s^{(i)}[n] \oplus s^{(i)}[n-d]
\label{eq:residual}
\end{equation}
%
If $d \in \{d_1,\ldots,d_S\}$, the residual $r^{(i)}_d[n]$ equals
$\bigoplus_{k \ne d} s^{(i)}[n-d_k] \oplus e[n]$; otherwise it contains
the full XOR combination plus noise.
The correlation statistic for offset $d$ is:
%
\begin{equation}
\rho(d) = \frac{1}{N \cdot T} \sum_{i=1}^{N} \sum_{n=W}^{255}
    \mathbf{1}\!\left[s^{(i)}[n] \oplus s^{(i)}[n-d] = 0\right]
\label{eq:rho}
\end{equation}

\subsection{Algorithm}

% TODO: fill in actual offset recovery steps and threshold details

\begin{enumerate}
  \item Compute $\rho(d)$ for all $d \in [1, 64]$.
  \item Threshold at $\hat{\rho} = $ \texttt{TODO} to select candidate offsets.
  \item Verify candidates by checking consistency across sequences.
  \item With recovered offsets $\hat{D}$, predict noise-free test bits by
        majority vote over $N$ sequences.
\end{enumerate}

\subsection{Results}

% TODO: fill in recovered offsets and prediction accuracy

Recovered offsets: $\hat{D} = \{$\texttt{TODO}$\}$.\\
Prediction accuracy on held-out noisy sequences: \texttt{TODO}\%.

% ══════════════════════════════════════════
% 3. APPROACH 2: GRADIENT DESCENT
% ══════════════════════════════════════════
\section{Approach 2: Soft Relaxation via Gradient Descent}
\label{sec:gd}

\subsection{Formulation}

We parametrize the offset selector as a soft weight vector
$\mathbf{w} \in \mathbb{R}^{64}$ with $w_d \in [0,1]$,
representing the probability that offset $d$ is active.
The soft predicted bit at position $n$ is:
%
\begin{equation}
\hat{s}[n] = \sigma\!\left(\sum_{d=1}^{64} w_d \cdot (2 s[n-d] - 1)\right)
\label{eq:soft}
\end{equation}
%
The training objective minimizes binary cross-entropy over all sequences
and positions $n \ge W_{\max}$:
%
\begin{equation}
\mathcal{L}(\mathbf{w}) = -\frac{1}{N \cdot T}
    \sum_{i,n} \Bigl[ s^{(i)}_n \log \hat{s}^{(i)}_n
    + (1-s^{(i)}_n)\log(1-\hat{s}^{(i)}_n) \Bigr]
\label{eq:loss}
\end{equation}
%
An $\ell_1$ sparsity penalty $\lambda \|\mathbf{w}\|_1$ encourages
selecting a small number of offsets.

\subsection{Optimization}

% TODO: fill in optimizer details, learning rate, epochs, lambda

\begin{itemize}
  \item Optimizer: Adam, lr $= $ \texttt{TODO}, $\lambda = $ \texttt{TODO}
  \item After convergence, threshold $w_d > 0.5$ to obtain discrete offsets.
  \item Binarized offsets are used for noise-free prediction.
\end{itemize}

\subsection{Results}

% TODO: fill in convergence curve and recovered offsets

Recovered offsets: $\hat{D} = \{$\texttt{TODO}$\}$.
Both methods agreed on the final offset set.

% ══════════════════════════════════════════
% 4. FAILED APPROACHES
% ══════════════════════════════════════════
\section{Failed Approaches}
\label{sec:failed}

\begin{table}[h]
\centering
\small
\begin{tabular}{@{}lp{4.6cm}@{}}
\toprule
Approach & Why it failed \\
\midrule
% TODO: fill in failed attempts
\texttt{TODO} & \texttt{TODO} \\
\texttt{TODO} & \texttt{TODO} \\
\bottomrule
\end{tabular}
\caption{Approaches that did not work.}
\label{tab:failed}
\end{table}

% ══════════════════════════════════════════
% 5. FINAL PREDICTION
% ══════════════════════════════════════════
\section{Final Prediction}
\label{sec:result}

Test prefix (64 bits, noise-free):
\begin{center}
\small\ttfamily
0000010100011010010101100101001110100011\\
110010110011010000111010
\end{center}

Predicted continuation (positions 64--255, 192 bits):
\begin{center}
\small\ttfamily
\texttt{TODO: 192-bit answer here}
\end{center}

Both Approach~1 and Approach~2 yielded identical offset sets,
giving high confidence in the prediction.

\vspace{4pt}
\hrule
\vspace{4pt}
{\footnotesize\color{gray}
SparseTap $\cdot$ KRAFTON AI R\&D Hackathon $\cdot$ Round 1, Day 2}

\end{document}
